% Modelo conceptual: entidades principales y sus relaciones (versión sin tipos
% ni claves de las tablas del Anexo D, figs/erd-bd.tex)
% Uso: % Modelo conceptual: entidades principales y sus relaciones (versión sin tipos
% ni claves de las tablas del Anexo D, figs/erd-bd.tex)
% Uso: % Modelo conceptual: entidades principales y sus relaciones (versión sin tipos
% ni claves de las tablas del Anexo D, figs/erd-bd.tex)
% Uso: % Modelo conceptual: entidades principales y sus relaciones (versión sin tipos
% ni claves de las tablas del Anexo D, figs/erd-bd.tex)
% Uso: \input{figs/modelo-conceptual} dentro de figure+\centering+\resizebox{\textwidth}{!}{...}

\begin{tikzpicture}[
  every node/.style={font=\small\sffamily},
  ent/.style={
    rectangle split, rectangle split parts=2,
    rectangle split part fill={aquaESI!70, white},
    draw=aquaESI!80, thick, rounded corners=3pt,
    align=left, text width=4.3cm
  },
  entref/.style={
    rectangle split, rectangle split parts=2,
    rectangle split part fill={gray!50, white},
    draw=gray!70, thick, rounded corners=3pt,
    align=left, text width=4.3cm
  },
  rel/.style={-, thick, black!70},
  card/.style={font=\small\sffamily\bfseries, fill=white, inner sep=1.5pt}
]

\node[ent] (usuario) at (0,0) {
  \textbf{Usuario}
  \nodepart{two}
  nombre, usuario, email\\
  contraseña (hash)\\
  idioma, confirmado
};

\node[ent] (analisis) at (8.6,0) {
  \textbf{Análisis}
  \nodepart{two}
  informe clínico enviado\\
  tarjetas de respuesta
};

\node[ent] (tarjeta) at (8.6,-3.6) {
  \textbf{Tarjeta de análisis}
  \nodepart{two}
  tipo: resumen o códigos
};

\node[ent] (codigo) at (0,-3.6) {
  \textbf{Código predicho}
  \nodepart{two}
  confianza\\
  estado: pendiente, validado\\
  o rechazado
};

\node[entref] (cie10) at (0,-7.0) {
  \textbf{Código CIE-10-ES}
  \nodepart{two}
  código, descripción\\
  metadatos clínicos
};

\draw[rel] (analisis.west) -- (usuario.east)
  node[card, pos=0.08, above]{N} node[card, pos=0.92, above]{1};
\draw[rel] (tarjeta.north) -- (analisis.south)
  node[card, pos=0.1, right]{N} node[card, pos=0.9, right]{1};
\draw[rel] (codigo.east) -- (tarjeta.west)
  node[card, pos=0.08, above]{N} node[card, pos=0.92, above]{1};
\draw[rel] (codigo.south) -- (cie10.north)
  node[card, pos=0.1, right]{N} node[card, pos=0.9, right]{1};

\end{tikzpicture}
 dentro de figure+\centering+\resizebox{\textwidth}{!}{...}

\begin{tikzpicture}[
  every node/.style={font=\small\sffamily},
  ent/.style={
    rectangle split, rectangle split parts=2,
    rectangle split part fill={aquaESI!70, white},
    draw=aquaESI!80, thick, rounded corners=3pt,
    align=left, text width=4.3cm
  },
  entref/.style={
    rectangle split, rectangle split parts=2,
    rectangle split part fill={gray!50, white},
    draw=gray!70, thick, rounded corners=3pt,
    align=left, text width=4.3cm
  },
  rel/.style={-, thick, black!70},
  card/.style={font=\small\sffamily\bfseries, fill=white, inner sep=1.5pt}
]

\node[ent] (usuario) at (0,0) {
  \textbf{Usuario}
  \nodepart{two}
  nombre, usuario, email\\
  contraseña (hash)\\
  idioma, confirmado
};

\node[ent] (analisis) at (8.6,0) {
  \textbf{Análisis}
  \nodepart{two}
  informe clínico enviado\\
  tarjetas de respuesta
};

\node[ent] (tarjeta) at (8.6,-3.6) {
  \textbf{Tarjeta de análisis}
  \nodepart{two}
  tipo: resumen o códigos
};

\node[ent] (codigo) at (0,-3.6) {
  \textbf{Código predicho}
  \nodepart{two}
  confianza\\
  estado: pendiente, validado\\
  o rechazado
};

\node[entref] (cie10) at (0,-7.0) {
  \textbf{Código CIE-10-ES}
  \nodepart{two}
  código, descripción\\
  metadatos clínicos
};

\draw[rel] (analisis.west) -- (usuario.east)
  node[card, pos=0.08, above]{N} node[card, pos=0.92, above]{1};
\draw[rel] (tarjeta.north) -- (analisis.south)
  node[card, pos=0.1, right]{N} node[card, pos=0.9, right]{1};
\draw[rel] (codigo.east) -- (tarjeta.west)
  node[card, pos=0.08, above]{N} node[card, pos=0.92, above]{1};
\draw[rel] (codigo.south) -- (cie10.north)
  node[card, pos=0.1, right]{N} node[card, pos=0.9, right]{1};

\end{tikzpicture}
 dentro de figure+\centering+\resizebox{\textwidth}{!}{...}

\begin{tikzpicture}[
  every node/.style={font=\small\sffamily},
  ent/.style={
    rectangle split, rectangle split parts=2,
    rectangle split part fill={aquaESI!70, white},
    draw=aquaESI!80, thick, rounded corners=3pt,
    align=left, text width=4.3cm
  },
  entref/.style={
    rectangle split, rectangle split parts=2,
    rectangle split part fill={gray!50, white},
    draw=gray!70, thick, rounded corners=3pt,
    align=left, text width=4.3cm
  },
  rel/.style={-, thick, black!70},
  card/.style={font=\small\sffamily\bfseries, fill=white, inner sep=1.5pt}
]

\node[ent] (usuario) at (0,0) {
  \textbf{Usuario}
  \nodepart{two}
  nombre, usuario, email\\
  contraseña (hash)\\
  idioma, confirmado
};

\node[ent] (analisis) at (8.6,0) {
  \textbf{Análisis}
  \nodepart{two}
  informe clínico enviado\\
  tarjetas de respuesta
};

\node[ent] (tarjeta) at (8.6,-3.6) {
  \textbf{Tarjeta de análisis}
  \nodepart{two}
  tipo: resumen o códigos
};

\node[ent] (codigo) at (0,-3.6) {
  \textbf{Código predicho}
  \nodepart{two}
  confianza\\
  estado: pendiente, validado\\
  o rechazado
};

\node[entref] (cie10) at (0,-7.0) {
  \textbf{Código CIE-10-ES}
  \nodepart{two}
  código, descripción\\
  metadatos clínicos
};

\draw[rel] (analisis.west) -- (usuario.east)
  node[card, pos=0.08, above]{N} node[card, pos=0.92, above]{1};
\draw[rel] (tarjeta.north) -- (analisis.south)
  node[card, pos=0.1, right]{N} node[card, pos=0.9, right]{1};
\draw[rel] (codigo.east) -- (tarjeta.west)
  node[card, pos=0.08, above]{N} node[card, pos=0.92, above]{1};
\draw[rel] (codigo.south) -- (cie10.north)
  node[card, pos=0.1, right]{N} node[card, pos=0.9, right]{1};

\end{tikzpicture}
 dentro de figure+\centering+\resizebox{\textwidth}{!}{...}

\begin{tikzpicture}[
  every node/.style={font=\small\sffamily},
  ent/.style={
    rectangle split, rectangle split parts=2,
    rectangle split part fill={aquaESI!70, white},
    draw=aquaESI!80, thick, rounded corners=3pt,
    align=left, text width=4.3cm
  },
  entref/.style={
    rectangle split, rectangle split parts=2,
    rectangle split part fill={gray!50, white},
    draw=gray!70, thick, rounded corners=3pt,
    align=left, text width=4.3cm
  },
  rel/.style={-, thick, black!70},
  card/.style={font=\small\sffamily\bfseries, fill=white, inner sep=1.5pt}
]

\node[ent] (usuario) at (0,0) {
  \textbf{Usuario}
  \nodepart{two}
  nombre, usuario, email\\
  contraseña (hash)\\
  idioma, confirmado
};

\node[ent] (analisis) at (8.6,0) {
  \textbf{Análisis}
  \nodepart{two}
  informe clínico enviado\\
  tarjetas de respuesta
};

\node[ent] (tarjeta) at (8.6,-3.6) {
  \textbf{Tarjeta de análisis}
  \nodepart{two}
  tipo: resumen o códigos
};

\node[ent] (codigo) at (0,-3.6) {
  \textbf{Código predicho}
  \nodepart{two}
  confianza\\
  estado: pendiente, validado\\
  o rechazado
};

\node[entref] (cie10) at (0,-7.0) {
  \textbf{Código CIE-10-ES}
  \nodepart{two}
  código, descripción\\
  metadatos clínicos
};

\draw[rel] (analisis.west) -- (usuario.east)
  node[card, pos=0.08, above]{N} node[card, pos=0.92, above]{1};
\draw[rel] (tarjeta.north) -- (analisis.south)
  node[card, pos=0.1, right]{N} node[card, pos=0.9, right]{1};
\draw[rel] (codigo.east) -- (tarjeta.west)
  node[card, pos=0.08, above]{N} node[card, pos=0.92, above]{1};
\draw[rel] (codigo.south) -- (cie10.north)
  node[card, pos=0.1, right]{N} node[card, pos=0.9, right]{1};

\end{tikzpicture}
